% !TeX root = mainstarkov.tex
% !TeX program = xelatex
% !TeX encoding = UTF-8
\section*{ВВЕДЕНИЕ}
\addcontentsline{toc}{section}{ВВЕДЕНИЕ}

Современные системы навигации в замкнутых пространствах требуют не только точного построения маршрутов, но и высокой производительности алгоритмов поиска пути, особенно в условиях динамически изменяющейся обстановки. В чрезвычайных ситуациях, когда определённые участки маршрута могут стать недоступными, система должна за миллисекунды перестроить оптимальный путь до ближайшего безопасного выхода.

Традиционные планы эвакуации решают только часть этой задачи. Они показывают общую схему движения, но остаются статичными: план не определяет местоположение конкретного пользователя, не учитывает заблокированный коридор и не перестраивает путь при изменении обстановки. При этом требования к эвакуационным путям и пожарной безопасности закреплены в нормативных документах, а значит программная система должна не подменять эти требования, а дополнять их прикладным механизмом построения маршрута.

В основе такой системы лежит представление здания в виде графа. Помещения, лестницы, лифты, выходы, QR-якоря и промежуточные точки маршрута рассматриваются как узлы, а допустимые переходы между ними – как рёбра. Это позволяет применять алгоритмы поиска кратчайшего пути, в том числе алгоритм Дейкстры, и учитывать дополнительные условия: смену этажа, недоступность отдельных участков, выбор ближайшего безопасного выхода. Теоретическая база подобных задач опирается на методы работы с графами и алгоритмы построения маршрутов.

Актуальность работы определяется необходимостью обеспечить надёжное и быстрое построение эвакуационных маршрутов в реальном времени. Алгоритмы маршрутизации должны обрабатывать графы с сотнями узлов и рёбер, учитывать межэтажные переходы, адаптироваться к изменениям в графе и гарантировать нахождение оптимального пути при наличии решения. Согласно критериям качества программного обеспечения по ГОСТ Р ИСО/МЭК 25010, в наибольшей степени проявляются характеристики функциональной пригодности (корректность маршрутизации), эффективности по производительности (время построения маршрута) и надёжности (устойчивость к блокировкам). Поэтому при реализации недостаточно только запрограммировать алгоритм: необходимо определить подходящие структуры данных, обеспечить обработку динамических ограничений и верифицировать корректность тестами.

\emph{Цель настоящей работы} -- реализовать и протестировать алгоритмы динамической маршрутизации для системы построения эвакуационных маршрутов, обеспечивающие построение оптимальных путей с учётом заблокированных участков и режима чрезвычайной ситуации.

Для достижения поставленной цели необходимо решить \emph{следующие задачи}:
\begin{itemize}
\item изучить существующие алгоритмы поиска кратчайших путей на графах;
\item спроектировать структуры данных для эффективного хранения и обработки навигационного графа;
\item реализовать алгоритм Дейкстры с поддержкой динамического исключения узлов;
\item реализовать алгоритм поиска ближайшего эвакуационного выхода с учётом текущего положения;
\item построить механизм формирования множества исключаемых узлов для перестроения маршрута;
\item провести модульное и системное тестирование алгоритмов на разных сценариях;
\item оценить быстродействие и правильность работы алгоритмов.
\end{itemize}

\emph{Структура и объём работы} -- Отчет состоит из введения, 4 разделов
основной части, заключения, списка использованных источников, приложений и графического материала на отдельных листах. Объём текста выпускной
квалификационной работы составляет \formbytotal{lastpage}{страниц}{у}{ы}{}.

\emph{Во введении} сформулирована цель работы, поставлены задачи разработки, описана структура работы, приведено краткое содержание каждого из разделов.

\emph{В первом разделе} на стадии анализа предметной области рассмотрены алгоритмы поиска кратчайших путей на графах, проведён их сравнительный анализ, рассмотрены способы представления графа в памяти и проанализированы существующие реализации систем indoor-навигации.

\emph{Во втором разделе} на стадии технического задания приводятся требования к разрабатываемой программной системе.

\emph{В третьем разделе} на стадии технического проектирования представлены проектные решения для программной системы.

\emph{В четвертом разделе} приводится список функций, классов и их методов, использованных при разработке программной системы, производится тестирование разработанной программной системы

В заключении излагаются основные результаты работы, полученные в ходе разработки.

В приложении А представлен графический материал.
В приложении Б представлены фрагменты исходного кода. 